\section{Dimensionality Reduction}
In Dimensionality Reduction, we assume that although the data is
high-dimensional, it lies on or near a low-dimensional subspace or manifold.
\subsection{Principal Component Analysis}
Considering a linear subspace $\mathbb{R}^{d}$ within $\mathbb{R}^{m}$ where
$d<<m$, we can reduce the dimensionality of our data, since we only require
$\mathbb{R}^{d}$ coordinates to describe the point. For this problem, we
consider Principal Component Analysis, a method of linear dimensionality
reduction.
\subsubsection{Problem Setting}
We assume our data consists $n$ instances: ${\left\{ \mathbf{x}^{(i)}
\right\}}_{i=1}^{n}$ where $\mathbf{x}^{(i)}\in \mathbb{R}^{m}$ of which we are
projecting (lienarly) onto a space $\mathbb{R}^{d}$ where $d<m$ such that the
variance of projected data is maximised. The space is spanned by a set of
orthonormal basis vectors ${\left\{ \mathbf{u}^{[i]} \right\}}^{d}_{i=1}$ where
$\mathbf{u}^{[i]}\in \mathbb{R}^{m}$.\\
Additionaly, we define the mean of the projected data onto each
basis vector $\mathbf{u}^{[j]}\cdot\bar{\mathbf{x}}$ where $\bar{\mathbf{x}}$ is the
sample mean:
\[\bar{\mathbf{x}}=\frac{1}{n}\sum_{i=1}^{n} \mathbf{x}^{(i)}.\]
\subsubsection{Projected Variance Maximisation}
We adopt the view of PCA of Projected Variance Maximisation, where the idea is
to project our high-dimensional data onto a lower-dimensional subspace such
that sample variance is maximally preserved. 
\begin{figure}[h!]
  \centering
  \includegraphics[width=0.3\textwidth]{variance-max}
\end{figure}
\begin{definition}[Sample Variance]
  The sample variance of the projected data is given by:
  \[\frac{1}{n}\sum_{i=1}^{n} {\left( \mathbf{u}^{[j]}\cdot\mathbf{x}^{(i)}-\mathbf{u}^{[j]}\cdot\bar{\mathbf{x}} \right)}^{2}=\mathbf{u}^{[j]T}\mathbf{Su}^{[j]}\] 
  The sample covariance matrix $\mathbf{S}$ is defned by:
\[\mathbf{S}=\frac{1}{n}\sum_{i=1}^{n} \left( \mathbf{x}^{(i)}-\bar{\mathbf{x}} \right) {\left( \mathbf{x}^{(i)}-\bar{\mathbf{x}} \right)}^{T}=\frac{1}{n}\mathbf{X}^{T}\mathbf{X}\] 
\end{definition}
where the centred design matrix $\mathbf{X}$ is defined:
\[\mathbf{X}=\begin{bmatrix} {\left( \mathbf{x}^{(1)}-\bar{\mathbf{x}}
\right)}^{T}\\ {\left( \mathbf{x}^{(2)}-\bar{\mathbf{x}}
\right)}^{T}\\\vdots\\{\left( \mathbf{x}^{(n)}-\bar{\mathbf{x}}
\right)}^{T}\end{bmatrix}.\] 
\begin{remark}
  PCA can also be viewed as a search for the $d$-dimensional subspace minmising
  the reconstruction error:
  \[\sum_{i=1}^{n} \left\lVert \left( \mathbf{x}^{(i)}-\bar{\mathbf{x}} \right) -\sum_{j=1}^{d} \left( \mathbf{u}^{[j]}\cdot\left( \mathbf{x}^{(i)}-\bar{\mathbf{x}} \right)  \right) \mathbf{u}^{[j]} \right\rVert_2^{2}.\]
  \begin{figure}[h!]
    \centering
    \includegraphics[width=0.3\textwidth]{residual-min}
  \end{figure}
\end{remark}

\subsubsection{Optimisation}
Re-writing our optimisation problem, we have:
\[\argmax_{{\left\{ \mathbf{u}^{[j]} \right\}}_{j=1}^{d}}L\quad\text{where:}\quad L=\sum_{j=1}^{d} \mathbf{u}^{[j]T}\mathbf{Su}^{[j]}.\]
for some orthonormal ${\left\{ \mathbf{u}^{[i]} \right\}}^{d}_{i=1}$.\\
In this algorithm, we attempt to learn each dimension greedily.
We first search for the direction of highest variance, and write the objective
as follows:
\[\mathcal{L}(\mathbf{u}^{[1]},\lambda^{[1]})=\mathbf{u}^{[1]T}\mathbf{Su}^{[1]}-\lambda^{[1]}\left( \mathbf{u}^{[1]}\cdot\mathbf{u}^{[1]}-1 \right)\] 
for some lagrange multiplier $\lambda^{[1]}$ associated with the orthogonality
constraint. Seeking stationarity w.r.t. $\mathbf{u}^{[1]}$, we have:
\begin{align*}
  2\mathbf{Su}^{[1]}-2\lambda^{[1]}\mathbf{u}^{[1]}=&0\\
  \implies\mathbf{Su}^{[1]}=&\lambda^{[1]}\mathbf{u}^{[1]}
\end{align*}
Hence $\mathbf{u}^{[1]}$ is an eigenvector of $\mathbf{S}$ associated with the
eigenvalue $\lambda^{[1]}$.  To find the variance of the projected data, we
have:
\[\mathbf{u}^{[1]T}\mathbf{Su}^{[1]}=\lambda^{[1]}\] 
To maximise this variance, we select the eigenvector associated with the
largest eigenvalue. We have found the first basis vector of greatest variance.
Now we seek to the second basis vector $\mathbf{u}^{[2]}$ to further increase
the projected variance, such that $\mathbf{u}^{[2]}\cdot \mathbf{u}^{[2]}=1$
and $\mathbf{u}^{[1]}\cdot \mathbf{u}^{[2]}=0$. We write the objective as
follows:
\[\mathcal{L}(\mathbf{u}^{[2]},\lambda^{[2]},\lambda^{[1][2]|})=\mathbf{u}^{[2]T}\mathbf{Su}^{[2]}-\lambda^{[2]}\left( \mathbf{u}^{[2]}\cdot\mathbf{u}^{[2]}-1 \right)-\lambda^{[1][2]}\left( \mathbf{u}^{[2]}\cdot\mathbf{u}^{[1]} -0\right) \] 
for some lagrange multipliers $\lambda^{[2]}, \lambda^{[1][2]}$. Seeking
stationarity w.r.t. $\mathbf{u}^{[2]}$, we have:
\[2\mathbf{Su}^{[2]}-2\lambda^{[2]}\mathbf{u}^{[2]}-\lambda^{[1][2]}\mathbf{u}^{[1]}=0.\] 
Left multiplying with $\mathbf{u}^{[1]T}$, $\lambda^{[1][2]}=0$.
Hence:
\[\mathbf{Su}^{[2]}=\lambda^{[2]}\mathbf{u}^{[2]}.\] 
Similarly, we want to maximise the variance so we select the eigenvector
associated with the largest remaining eigenvalue. The solution proceeds in a
similar step-wise fashion to give:
\[{\left\{ \mathbf{u}^{[j]*}=\mathbf{q}_j \right\}}_{j=1}^{d}\]
implying a total projected variance of:
\[L^{*}=\sum_{j=1}^{d} \mathbf{q}_j^{T}\mathbf{Sq}_j=\sum_{j=1}^{d} \lambda_j.\] 
This solution assumes orthogonality of basis vectors, normality of different
basis vectors and the orientation of the first basis vector. Other choices are
also possible yielding similar solutions. While the original problem is
non-convex, it is possible to show that the greedy local optima gives rise to
\textbf{globally optimal points}.\\
PCA allows estimation of the input space as the space spanned by the first $d$
eigenvectors with the largest eigenvalues of the sample covariance matrix.
\begin{definition}[Ambient and Intrinsic Dimensionalities]
  The dimensionality of the input space is known as the ambient dimensionality
  of the data, and the dimensionality of the subspace upon which the data is
  assumed to lie is known as the intrinsic dimensionality of the data.
\end{definition}
\subsubsection{Pattern Stabilility}
To be sure that the training sample has discerned a reliable subspace
(intuitively, how extensible the subspace is to other data) we consider two
other approaches:
\begin{itemize}
  \item \textbf{Generative Modelling}: Seeking a Gaussian latent variable model to be
    learned using MLE or Bayesian treatment
  \item \textbf{PAC Approach}: Bounding generalisation error resulting from
    projection $E_\mathcal{D}\left[ \mathbf{x}-\sum_{j=1}^{d} \left(
    \mathbf{u}^{[j]}\cdot\mathbf{x} \right) \mathbf{u}^{[j]} \right]$ with high
    probability. Bound consists terms in sample residual eigenspectrum and
    complexity of data space. With a subspace capturing a high proportion of
    the variance with a small relative dimensionality, we have low
    generalisation error.
\end{itemize}

\subsection{Probabilistic PCA}
In Probabilistic PCA, we seek a Gaussian latent variable model and learn its parameters
as a probabilistic generalisation of PCA.\@
\subsubsection{PPCA Model}
We assume that each point has an unknown latent variable $\mathbf{z}\in
\mathbb{R}^{d}$ associated with it corresponding to its position in the
principal component subspace. We assume this variable is a outcome of a Gaussian random variable
$\mathcal{Z}$:
\[\mathbf{z}\sim\mathcal{N}\left( \mathbf{0},\mathbf{I}_d \right).\] 
Additionally, each $\mathbf{x}$ is a outcome of a Gaussian random variable $\mathcal{X}$:
\[\mathbf{x|z}\sim\mathcal{N}\left( \mathbf{Wz+\mu},\sigma^{2}\mathbf{I}_m \right).\]
where $\mathbf{W}\in \mathbb{R}^{m\times d}$ defining the directions of the
principal subspace, $\mu\in \mathbb{R}^{m},\sigma\in \mathbb{R}^{+}$.\\

Intuitively, we can understand the model as follows:
\begin{itemize}
  \item We draw an outcome for the latent variable $\mathbf{z}$.
  \item The observed variable is sample conditional on the latent variable $\mathbf{x}|\mathbf{z}$.
  \item We observe residual noise $\varepsilon$ such that:
    \[\mathbf{x}=\mathbf{Wz+\mu+\varepsilon}.\] 
    \[\mathbf{z}\sim\mathcal{N}\left( \mathbf{0},\mathbf{I}_d \right).\] 
    \[\mathbf{\varepsilon}\sim\mathcal{N}\left( \mathbf{0},\sigma^{2}\mathbf{I}_m \right).\] 
    \begin{remark}
      Given a marginal distribution for $\mathbf{\tilde{x}}$ and a conditional
      Gaussian distribution for $\mathbf{\tilde{y}}$ given
      $\mathbf{\tilde{x}}$:
      \[\mathbf{\tilde{x}}\sim\mathcal{N}(\mathbf{\mu},\mathbf{\Lambda}^{-1})\] 
      \[\mathbf{\tilde{y}|\tilde{x}}\sim\mathcal{N}(\mathbf{A\tilde{x}+b},\mathbf{L}^{-1})\] 
      for some $\mathbf{\tilde{x}}\in \mathbb{R}^{n},\mathbf{\tilde{y}}\in \mathbb{R}^{m},\mathbf{\mu}\in \mathbb{R}^{n},\mathbf{\Lambda}\in \mathbb{R}^{n\times n},\mathbf{A}\in \mathbb{R}^{m\times n},\mathbf{b}\in \mathbb{R}^{m},\mathbf{L}\in \mathbb{R}\in ^{m\times m}$. Then:
    \[\mathbf{\tilde{y}}\sim\mathcal{N}\left( \mathbf{A\mu+b,L}^{-1}+\mathbf{A\Lambda}^{-1}\mathbf{A^{T}} \right).\] 
      \[\mathbf{\tilde{x}|\tilde{y}}\sim\mathcal{N}(\mathbf{\Sigma} \left[ \mathbf{A}^{T}\mathbf{L(\tilde{y}-b)+\Lambda\mu} \right],\mathbf{\Sigma})\] 
      where $\Sigma={(\mathbf{\Lambda}+\mathbf{A}^{T}\mathbf{LA})}^{-1}$
    \end{remark}
    Given this, we have:
    \[\mathbf{x}\sim\mathcal{N}\left( \mathbf{\mu},\mathbf{WW}^{T}+\sigma^{2}\mathbf{I}_m \right).\] 
\end{itemize}
Since the covariance matrix of the distribution of $\mathbf{x}$ is symmetric,
we can interpret $p_\mathcal{X}(\mathbf{x})$ as a density determined by
$\sigma^{2}$ of taking an isotropic Gaussian `spray can' across the principal
subspace. We have a `pancake' shaped distribution.
\subsubsection{Rotational Invariance}
Suppose $\mathbf{R}$ an orthogonal (rotation) matrix (hence $\mathbf{R}^{T}\mathbf{R}=\mathbf{I}$).
Applying it to the latent space coordinate matrix $\mathbf{W}$:
\begin{align*}
  \mathbf{\tilde{W}}=&\mathbf{WR}\\
  \implies\mathbf{\tilde{W}\tilde{W}}^{T}=&\mathbf{WRR}^{T}\mathbf{W}^{T}\\
  =&\mathbf{WW}^{T}
\end{align*}
The distribution of $\mathbf{x}$ can be characterised by any
$\mathbf{\tilde{W}}$. THis is the analogue of the non-uniqueness in PCA.\@
\subsubsection{Optimisation}
We find the log-likelihood function to optimise as follows:
\begin{align*}
  \mathcal{L}=&\sum_{i=1}^{n} \ln p_\mathcal{X}\left( \mathbf{x}^{(i)};\mathbf{W,\mu},\sigma^{2} \right) \\
  =&-\frac{nm}{2}\ln\left( 2\pi \right) -\frac{n}{2}\ln\left| \mathbf{C} \right| -\frac{1}{2}\sum_{i=1}^{n} {\left( \mathbf{x}^{(i)}-\mathbf{\mu} \right)}^{T}\mathbf{C}^{-1}\left( \mathbf{x}^{(i)}-\mathbf{\mu} \right) 
\end{align*}
It can be demonstrated that the following closed form solutions flow from the optimisation of this
function (Tipping \& Bishop, 1999):
\[\mathbf{\mu}_{MLE}=\mathbf{\bar{x}}.\]
\[\sigma_{MLE}=\frac{1}{m-d}\sum_{i=d+1}^{m} \lambda_i.\] 
\[\mathbf{W}_{MLE}=\mathbf{Q}{\left( \mathbf{\Lambda}-\sigma^{2}\mathbf{I} \right)}^{\frac{1}{2}\mathbf{R}}.\] 
where $\mathbf{Q}\in \mathbb{R}^{m\times d}$ whose columns are given by the
leading $d$ eigenvectors of the covariance matrix $\mathbf{S}$,
$\Lambda=\diag(\lambda_1,\lambda_2,\ldots,\lambda_d)$ contains the $d$
eigenvalues associated with these eigenvectors, and $\mathbf{R}$ is an
arbitrary orthogonal matrix. We can set $\mathbf{R=I}$ without loss of
generality, when $\mathbf{W}$ are the principal component eigenvectors scaled
by the square root of the variance parameters ${(\lambda_i-\sigma^{2})}^{\frac{1}{2}}$.
Hence, the variance of $p_\mathcal{X}(\mathbf{x})$ in the $\mathbf{q}_i$
direction is $\lambda_i$.

